\begin{table}[htbp]
  \centering
  \small
  \caption{Emojis Más Frecuentes en el Cuerpo del Texto y su Categorización Semántica}
  \label{tab:top_emojis}
  \begin{tabular}{llrr}
    \toprule
    Emoji / Descripción & Emoción Asignada & Frecuencia ($n$) & Proporción (\%) \\
    \midrule
    \texttt{[U+1F44D]} Pulgar arriba (\emph{thumbs up})   & Aprobación / Gratitud  & 412 & 23,8 \\
    \texttt{[U+1F602]} Risa con lágrimas (\emph{joy})     & Diversión / Alegría    & 385 & 22,2 \\
    \texttt{[U+1F64F]} Manos juntas (\emph{pray})         & Gratitud / Ruego        & 310 & 17,9 \\
    \texttt{[U+1F914]} Cara pensativa (\emph{thinking})   & Duda / Anticipación     & 280 & 16,2 \\
    \texttt{[U+1F525]} Fuego (\emph{fire})                & Entusiasmo / Intensidad & 195 & 11,3 \\
    \texttt{[U+1F62D]} Llanto intenso (\emph{sob})        & Tristeza / Frustración  & 140 & 8,1  \\
    \bottomrule
  \end{tabular}
  \vspace{4pt}
  \begin{minipage}{0.92\textwidth}
    \footnotesize
    \textit{Nota.} Extraídos del campo \texttt{emoji\_list} en la tabla \texttt{mensajes}. No se purgan durante la tokenización para conservar su función como operadores léxicos afectivos independientes.
  \end{minipage}
\end{table}
